% Generated from locale/tr/content/first-order-logic/completeness/completeness-thm.tex; do not hand-edit.


\iftag{FOL}
      {\olfileid[tr]{fol}{com}{cth}}
      {\olfileid[tr]{pl}{com}{cth}}

\olsection{Tamlık Teoremi}

\begin{explain}
Sonuçlarımızı bir araya getirdiğimizde tamlık teoremine ulaşırız.
\end{explain}

\begin{thm}[Tamlık Teoremi]
\ollabel{thm:completeness}
$\Gamma$ bir !!{sentence}s kümesi olsun. $\Gamma$ tutarlıysa
sağlanabilirdir.
\end{thm}

\begin{proof}
$\Gamma$'nın tutarlı olduğunu varsayalım.\iftag{FOL}{
  \olref[hen]{lem:henkin} uyarınca $\Gamma'\supseteq\Gamma$ olacak
  biçimde doygun ve tutarlı bir~$\Gamma'$ kümesi vardır.}{}
\olref[lin]{lem:lindenbaum} uyarınca
$\Gamma^*\supseteq\iftag{FOL}{\Gamma'}{\Gamma}$ olacak biçimde tutarlı
ve !!{complete} bir~$\Gamma^*$ kümesi vardır.
\iftag{FOL}{
  $\Gamma'\subseteq\Gamma^*$ olduğundan her $!A(x)$ !!{formula}
  ifadesi için $\Gamma^*$,
  \iftag{prvEx}
        {$\lexists[x][!A(x)]\lif !A(c)$}
        {$\lnot\lforall[x][!A(x)]\lif\lnot !A(c)$}
  biçiminde !!a{sentence} içerir ve dolayısıyla $\Gamma^*$ doygundur.
  $\Gamma$,~$\eq$ içermiyorsa}{ }
\olref[mod]{lem:truth} uyarınca
$\iftag{FOL}{\Sat{M(\Gamma^*)}}{\pSat{v(\Gamma^*)}}{!A}$ ancak ve ancak
$!A\in\Gamma^*$ ise. Özellikle her $!A\in\Gamma$ için
$\iftag{FOL}{\Sat{M(\Gamma^*)}}{\pSat{v(\Gamma^*)}}{!A}$ olur;
dolayısıyla $\Gamma$ sağlanabilirdir.\iftag{FOL}{
  $\Gamma$,~$\eq$ içeriyorsa \olref[ide]{lem:truth} uyarınca bütün
  !!{sentence}s~$!A$ için $\Sat{\equivclass{M}{\approx}}{!A}$ ancak ve
  ancak $!A\in\Gamma^*$ ise. Özellikle her $!A\in\Gamma$ için
  $\Sat{\equivclass{M}{\approx}}{!A}$ olur; dolayısıyla $\Gamma$
  sağlanabilirdir.}{}
\end{proof}

\begin{cor}[Tamlık Teoremi, İkinci Biçim]
\ollabel{cor:completeness}
Her $\Gamma$ ve her !!{sentence}s~$!A$ için: $\Gamma\Entails !A$ ise
$\Gamma\Proves !A$.
\end{cor}

\begin{proof}
\olref{cor:completeness} ile \olref{thm:completeness} içindeki
$\Gamma$ değişkenlerinin evrensel niceleyicili olduğuna dikkat edin.
Karışıklığı önlemek için \olref{thm:completeness} sonucunu farklı bir
değişkenle yeniden ifade edelim: herhangi bir !!{sentence}s
kümesi~$\Delta$ için, $\Delta$ tutarlıysa sağlanabilirdir. Karşıt-tersi
gereği $\Delta$ sağlanamazsa tutarsızdır. Sonucu ispatlamak için
bunu kullanacağız.

$\Gamma\Entails !A$ olduğunu varsayalım.
\olref[syn][sem]{prop:entails-unsat} uyarınca
$\Gamma\cup\{\lnot !A\}$ sağlanamaz. $\Gamma\cup\{\lnot !A\}$
kümesini~$\Delta$ olarak alırsak \olref{thm:completeness} teoreminin
önceki biçimi $\Gamma\cup\{\lnot !A\}$ kümesinin tutarsız olduğunu verir.
\iftag{FOL}{%
  \tagrefs{prfAX/{fol:axd:prv:prop:prov-incons},
    prfSC/{fol:seq:prv:prop:prov-incons},
    prfND/{fol:ntd:prv:prop:prov-incons},
    prfTab/{fol:tab:prv:prop:prov-incons}}}{%
  \tagrefs{prfAX/{pl:axd:prv:prop:prov-incons},
    prfSC/{pl:seq:prv:prop:prov-incons},
    prfND/{pl:ntd:prv:prop:prov-incons},
    prfTab/{pl:tab:prv:prop:prov-incons}}}
uyarınca $\Gamma\Proves !A$.
\end{proof}

\tagprob{FOL}
\begin{prob}
\olref[fol][com][cth]{cor:completeness} sonucunu kullanarak
\olref[fol][com][cth]{thm:completeness} teoremini ispatlayın ve böylece
tamlık teoreminin iki biçiminin denk olduğunu gösterin.
\end{prob}
\tagendprob

\tagprob{notFOL}
\begin{prob}
\olref[pl][com][cth]{cor:completeness} sonucunu kullanarak
\olref[pl][com][cth]{thm:completeness} teoremini ispatlayın ve böylece
tamlık teoreminin iki biçiminin denk olduğunu gösterin.
\end{prob}
\tagendprob

\tagprob{FOL}
\begin{prob}
!!a{derivation} sisteminin tam olabilmesi için kuralları, sağlanamaz
her kümenin tutarsız olduğunu ispatlayacak kadar güçlü olmalıdır.
Tamlığı ispatlamak için !!{derivation} kurallarından hangileri
gereklidir? Bu kurallardan herhangi biri ispatın hiçbir yerinde
kullanılmamış mıdır? Soruları yanıtlamak için
\olref[fol][com][cth]{thm:completeness} teoreminin ispatına götüren
sonuçların hangisinde hangi !!{derivation} kurallarının kullanıldığını
gösteren bir liste ya da şema hazırlayın. Bu ispatlarda kuralların
örtük kullanımlarını da belirtin.
\end{prob}
\tagendprob

\tagprob{notFOL}
\begin{prob}
!!a{derivation} sisteminin tam olabilmesi için kuralları, sağlanamaz
her kümenin tutarsız olduğunu ispatlayacak kadar güçlü olmalıdır.
Tamlığı ispatlamak için !!{derivation} kurallarından hangileri
gereklidir? Bu kurallardan herhangi biri ispatın hiçbir yerinde
kullanılmamış mıdır? Soruları yanıtlamak için
\olref[pl][com][cth]{thm:completeness} teoreminin ispatına götüren
sonuçların hangisinde hangi !!{derivation} kurallarının kullanıldığını
gösteren bir liste ya da şema hazırlayın. Bu ispatlarda kuralların
örtük kullanımlarını da belirtin.
\end{prob}
\tagendprob

